{{دیوان
| عنوان = فرمانِ خدا
| ماخذ = ویکی ماخذ
| ماخذ_لنک = https://ur.wikisource.org/wiki/%D9%81%D8%B1%D9%85%D8%A7%D9%86%D9%90_%D8%AE%D8%AF%D8%A7
| مقابلہ = اقبال سائبر لائبریری (اقبال اکادمی پاکستان)
| مقابلہ_لنک = https://www.iqbalcyberlibrary.net
}}
=== (فرشتوں سے) ===

<poem>
اُٹھّو! مری دنیا کے غریبوں کو جگا دو
کاخِ اُمَرا کے در و دیوار ہِلا دو

گرماؤ غلاموں کا لہُو سوزِ یقیں سے
کُنجشکِ فرومایہ کو شاہیں سے لڑا دو

سُلطانیِ جمہور کا آتا ہے زمانہ
جو نقشِ کُہَن تم کو نظر آئے، مِٹا دو

جس کھیت سے دہقاں کو میسرّ نہیں روزی
اُس کھیت کے ہر خوشۂ گندم کو جلا دو

کیوں خالق و مخلوق میں حائل رہیں پردے
پِیرانِ کلیسا کو کلیسا سے اُٹھا دو

حق را بسجودے، صنَماں را بطوافے
بہتر ہے چراغِ حَرم و دَیر بُجھا دو

میں ناخوش و بیزار ہُوں مَرمَر کی سِلوں سے
میرے لیے مٹّی کا حرم اور بنا دو

تہذیبِ نوی کارگہِ شیشہ گراں ہے
آدابِ جُنوں شاعرِ مشرِق کو سِکھا دو!
</poem>
